% STATUS: DRAFT
\chapter{The punchline and how to read this book}\label{ch:punchline}

\begin{physmotiv}
Before any mathematics, here is the whole story in one page. Chapters 3--25 are the machinery needed to make each sentence below precise; Chapters 30--40 are the physics it is aimed at. If you are impatient, read this chapter first, then use the reading paths in \cref{sec:paths}.
\end{physmotiv}

\section{The story in four acts}

\textbf{Act I: Quantum field theory has no entropy.} The algebra $\A(R)$ of observables in a spacetime region $R$ of a QFT is a von Neumann algebra of \emph{type III$_1$}. Such an algebra has no trace and no density matrices, and consequently no von Neumann entropy $-\Tr\rho\log\rho$. The ``entropy'' one computes on a lattice is a cutoff-dependent artifact (\cref{ch:why_algebras,ch:typeIII_local}).

\textbf{Act II: Modular theory is hidden dynamics.} Any state $\Omega$ that is cyclic and separating for an algebra $\M$ determines a one-parameter group $\sigma_t$ of automorphisms of $\M$, the \emph{modular flow}. If $\M$ were $\BH$ with density matrix $\rho$, then $\sigma_t(a)=\rho^{it}a\rho^{-it}$, so the generator is $-\log\rho$, the \emph{modular Hamiltonian}. For a type III algebra there is no $\rho$, but the modular flow still exists (Tomita--Takesaki, \cref{ch:tomita}) and is not inner. For a Rindler wedge in the Minkowski vacuum, it is the boost (\cref{ch:rs_bw}).

\textbf{Act III: Crossed products make time inner.} Given $\M$ and its modular flow, the \emph{crossed product} $\M\rtimes_\sigma\R$ is the algebra generated by $\M$ together with a new operator that implements the flow. The theorem of Connes and Takesaki is that if $\M$ is type III$_1$, then the crossed product is type II$_\infty$: it \emph{has} a trace, and so entropy is defined up to a state-independent additive ambiguity (\cref{ch:crossed,ch:takesaki,ch:entropy_II}).

\textbf{Act IV: Gravity does it automatically.} In gravity, the Hamiltonian is a constraint, and its generator is a gauge redundancy. When one includes an observer with energy $q$ in de Sitter space and imposes the constraint, the algebra of gauge-invariant observables of the static patch is exactly such a crossed product (\cref{ch:clpw}). If the observer's energy is bounded below, the trace is finite and the algebra is of type II$_1$: there is a maximum entropy state, the de Sitter vacuum, and the entropy of semiclassical states is the \emph{generalized entropy} $\Sgen=A/4G+S_{\rm bulk}$ (\cref{ch:anatomy,ch:gen_entropy}).

\begin{center}
\begin{tikzpicture}[node distance=2.2cm, every node/.style={font=\small}]
\node[draw,rounded corners,align=center,fill=physcol!8] (a) {QFT region\\ type III$_1$};
\node[draw,rounded corners,align=center,right=of a,fill=physcol!8] (b) {+ modular flow\\ $\sigma_t$ (outer)};
\node[draw,rounded corners,align=center,right=of b,fill=physcol!8] (c) {crossed product\\ $\M\rtimes\R$};
\node[draw,rounded corners,align=center,right=of c,fill=physcol!8] (d) {type II\\ trace, entropy};
\draw[-{Stealth}] (a)--(b); \draw[-{Stealth}] (b)--(c); \draw[-{Stealth}] (c)--(d);
\node[below=0.8cm of c,align=center,font=\small\itshape] (g) {gravitational constraint $\;\Rightarrow\;$ this is automatic};
\draw[dashed,-{Stealth}] (g.north) -- (c.south);
\end{tikzpicture}
\end{center}

\subsection*{A table of types and entropy}
\begin{center}
\renewcommand{\arraystretch}{1.25}
\begin{tabular}{@{}llll@{}}
\toprule
Type & Example & Trace? & Entropy\\
\midrule
I$_n$ & $n\times n$ matrices & finite & yes\\
I$_\infty$ & $\BH$, $\dim\HH=\infty$ & infinite (not on identity) & yes, for trace-class $\rho$\\
II$_1$ & hyperfinite II$_1$ factor & finite, $\Tr\id=1$ & yes, bounded above\\
II$_\infty$ & II$_1\,\bar\otimes\,$I$_\infty$ & infinite & yes, up to shift\\
III$_\lambda$ & local QFT algebras (III$_1$) & none & none\\
\bottomrule
\end{tabular}
\end{center}
This table is the mnemonic for the first half of the book: Part III builds the vocabulary, Part IV explains why III exists, and Part V explains how to leave it.

\section{Dependency map}

\begin{center}
\begin{tikzpicture}[node distance=0.9cm and 1.1cm, every node/.style={draw,rounded corners,font=\small,minimum width=2.6cm,align=center}]
\node (I) {I Toolkit\\ Ch.\ 3--7};
\node[right=of I] (II) {II \cstar, states, GNS\\ Ch.\ 8--12};
\node[right=of II] (III) {III vN algebras\\ Ch.\ 13--17};
\node[below=of III] (IV) {IV KMS, modular\\ Ch.\ 18--21};
\node[left=of IV] (V) {V Crossed prod.\\ Ch.\ 22--25};
\node[below=of V] (VII) {VII Gravity, dS\\ Ch.\ 30--40};
\node[right=of VII] (VI) {VI AQFT\\ Ch.\ 26--29};
\node[below=of VI,xshift=-1.8cm] (VIII) {VIII BH/holography\\ Ch.\ 41--44};
\node[right=of VI] (IX) {IX NCG\\ Ch.\ 45--52};
\draw[-{Stealth}] (I)--(II); \draw[-{Stealth}] (II)--(III); \draw[-{Stealth}] (III)--(IV);
\draw[-{Stealth}] (IV)--(V); \draw[-{Stealth}] (V)--(VII); \draw[-{Stealth}] (IV)--(VI);
\draw[-{Stealth}] (VI)--(VII); \draw[-{Stealth}] (VII)--(VIII); \draw[-{Stealth}] (IV)--(IX);
\end{tikzpicture}
\end{center}

\section{Reading paths}\label{sec:paths}

\begin{description}[leftmargin=1em]
\item[Path A (de Sitter fast track).] All chapters labelled ``core'': Parts 0--VII minus Chapters 6 and 29. Roughly $35$ chapters. You will have all the mathematics needed for \cref{ch:clpw}.
\item[Path B (physicist in a hurry).] Chapters 1--2, 3--4 (skim), 8--10, 13--16, 18--20, 22--25, then 32--36. This leaves out inequivalent representations and AQFT; accept the statement that local algebras are type III$_1$.
\item[Path C (mathematically curious).] Everything, including the boxes marked ``skippable'' and Part IX.
\item[Path D (noncommutative geometry).] Chapters 1--2, 3--5, 8, 13--16, then Part IX (which has its own appendices).
\end{description}

\section{Conventions}

Chapters follow a fixed template: a physics-motivation box, definitions and examples, a ``what could go wrong'' box, a key-idea summary, and exercises. Each Part ends with a dictionary translating the mathematical vocabulary back into physics. We use $\hbar=c=k_B=1$. Our conventions for modular theory are fixed once and for all in \cref{ch:tomita} and summarized in the style guide.

\begin{unsettled}
Chapters 37--40 and 52 describe results from 2022 onward that are still evolving; the status of these chapters is flagged in the text, and citations to the recent literature should be checked against current versions of the papers.
\end{unsettled}

\section*{Exercises}
\begin{exercise}
For which of the types in the table above does the identity operator have a finite trace? What does that imply about the existence of a maximum entropy state?
\end{exercise}
\begin{exercise}
Let $\rho$ be a density matrix on $\C^n$ and $\M=M_n(\C)$. Show that $\sigma_t(a)=\rho^{it}a\rho^{-it}$ is an inner automorphism, i.e.\ implemented by a unitary in $\M$. Why does inner-ness make the crossed product trivial in this case? (Hint: it will be isomorphic to $\M\otimes\Linf(\R)$.)
\end{exercise}
